\chapter{Hysteroscopy}

\section*{Overview}
Hysteroscopy is the inspection of the uterine cavity by endoscopy. A hysteroscope is inserted through the cervix to diagnose and treat intrauterine pathology.

\section*{Alternative Names}
Uterine endoscopy

\section*{Principle}
A thin telescope with a light and camera is passed through the cervix into the uterine cavity, which is distended with fluid or gas. The endometrial cavity is inspected directly, and instruments allow biopsy, removal of polyps or fibroids, and treatment of adhesions or septa.
\section*{History}
Pantaleoni performed the first hysteroscopy in 1869 using a modified cystoscope to identify a uterine polyp.

\section*{Why It Is Done}
Ordered for abnormal uterine bleeding, repeated miscarriages, infertility, checking for uterine fibroids or polyps, or locating a displaced IUD.

\section*{Is This a Primary or Secondary Test or Procedure}
Primary diagnostic and therapeutic tool for endometrial and intrauterine pathology.

\section*{How It Is Performed}
Performed under local, regional, or general anesthesia. The hysteroscope is passed through the vagina and cervix into the uterus. Liquid or gas is used to expand the cavity.

\section*{Preparation}
The procedure is usually scheduled when you are not menstruating. Fasting is required if regional or general anesthesia is planned.

\section*{Contraindications and Precautions}
Active pelvic infection, pregnancy, or known cervical or uterine cancer.

\section*{Risks}
Uterine perforation (<1%), cervical laceration, bleeding, pelvic infection, or fluid overload from the distending medium.

\section*{Aftercare and Recovery}
You may experience mild cramping and light vaginal bleeding or spotting for a few days. Avoid sexual intercourse until bleeding stops.

\section*{Results and Interpretation}
Enables direct visualization of the endometrium, fallopian tube openings, and structural anomalies. Biopsies or polypectomy can be performed.

\section*{Expected Results}
Normal cervix and uterine cavity; smooth endometrial lining; patent tubal ostia; no polyps or fibroids.

\section*{Follow-up}
Endometrial biopsy review, hormonal therapy, or surgical scheduling.

\section*{References}
\begin{enumerate}
  \item MedlinePlus. Hysteroscopy. U.S. National Library of Medicine; 2025.
  \item Current Medical Diagnosis and Treatment. McGraw-Hill; 2025.
\end{enumerate}

\clearpage
